% JUMP TO LINE 60, 75
\documentclass[preview, margin=0.6in]{standalone}
\usepackage[letterpaper,portrait,top=0.4in, left=0.6in, right=0.6in, bottom=1in]{geometry}

\usepackage{amsmath, amsfonts, amsthm, amssymb}
\usepackage{graphicx, float}
\usepackage{mathtools}
\usepackage{titlesec}
\usepackage{interval}
\usepackage{hyperref}
\usepackage{siunitx}
\usepackage{titling}
\usepackage{vwcol}
\usepackage{setspace}
\usepackage{empheq}
\usepackage{cancel}
\usepackage{esdiff}
\usepackage{multicol}
\usepackage{mdframed}
\usepackage{esdiff}
\usepackage{tikzsymbols}
\usepackage{multicol}
\usepackage{tikz}
\usepackage{varwidth}
\usepackage{parskip}
\usepackage{pgfplots}
\usepackage{mlmodern}
\pgfplotsset{compat=1.18}
\intervalconfig {
	soft open fences
}

\newcommand{\alignedintertext}[1]{%
  \noalign{%
    \vtop{\hsize=\linewidth#1\par
    \expandafter}%
    \expandafter\prevdepth\the\prevdepth
  }%
}

\newtheorem{lemma}{Lemma}

\newcommand*{\problem}[1]{\section*{Problem #1}}
\newcommand*{\aps}{\section*{AP Corner}}
\newcommand*{\deriv}[1][x]{\ensuremath{\dfrac{\mathrm{d}}{\mathrm{d}#1}}}
\newcommand*{\floor}[1]{\ensuremath{\lfloor #1\rfloor}}
\newcommand*{\lheqzero}{\ensuremath{\underset{\text{L'H}}{\overset{\left[\frac00\right]}{=}}}}
\newcommand*{\lheqinfty}{\ensuremath{\underset{\text{L'H}}{\overset{\left[\frac{\infty}{\infty}\right]}{=}}}}

\DeclareMathOperator{\DNE}{DNE}
\DeclareMathOperator{\sgn}{sgn}

\DeclareMathOperator{\arccsc}{arccsc}
\DeclareMathOperator{\arcsec}{arcsec}
\DeclareMathOperator{\arccot}{arccot}

%opening

\title{\vspace*{-40pt}Homework Set B.1}
\author{Jayden Li}
\date{\today}
% \allowdisplaybreaks
\postdisplaypenalty=100000

\begin{document}
\setstretch{1.25}
\fontsize{12pt}{12pt}\selectfont
\setlength{\abovedisplayskip}{\abovedisplayskip/2}
\setlength{\belowdisplayskip}{\belowdisplayskip/2}
\setlength{\parindent}{0pt}
\setlength{\parskip}{2ex plus 0.5ex minus 0.2ex}
\maketitle

\problem{1}
Let $v$ be the vector to the corner at $(1,1)$, so $v=(1,1)$. The $x$-axis is a line spanned by $\hat i=(1,0)$. Let $\theta\in\mathbb R$ be the angle between vectors $x$ and $\hat i$. By definition of the dot product:
\begin{equation*}
    v\cdot \hat i
	=1
	=\lVert v \rVert \lVert \hat i \rVert \cos\theta
	=\sqrt{2}(1)\cos\theta
	\implies \cos\theta=\frac{1}{\sqrt2}
	=\frac{\sqrt{2}}{2}
	\implies \theta=\frac{\pi}{4}=\boxed{45^\circ}
\end{equation*}

\problem{2}
Let $v$ be the vector from the origin to the corner, so $v=(1,1,1)$. The $x$-axis is spanned by $\hat i=(1,0,0)$.
\begin{equation*}
    v\cdot \hat i
	=1
	=\lVert v \rVert \lVert \hat i \rVert \cos\theta
	=\sqrt{3}(1)\cos\theta
	\implies \cos\theta=\frac{1}{\sqrt3}
	\implies \theta=\arccos\left(\frac{1}{\sqrt{3}}\right)
	\approx \boxed{54.7^\circ}
\end{equation*}

\problem{3}
$v=(1,1,1,1),\hat i=(1,0,0,0)$.
\begin{equation*}
    v\cdot \hat i
	=1
	=\lVert v \rVert \lVert \hat i \rVert \cos\theta
	=\sqrt{4}(1)\cos\theta
	\implies \cos\theta=\frac{1}{2}
	\implies \theta=\frac{\pi}{3}=\boxed{60^\circ}
\end{equation*}

\problem{4}
In $\mathbb R^n$, $v\cdot \hat i=1$, $\lVert \hat i \rVert =1$, and the magnitude of $v$ is:
\begin{equation*}
    \lVert v \rVert 
	=\sqrt{\sum_{i=1}^{n}v_i^2}
	=\sqrt{\sum_{i=1}^{n}1}
	=\sqrt{n}
\end{equation*}
Let $\theta$ be the angle between $v$ and $\hat i$. By definition of the dot product:
\begin{equation*}
	v\cdot\hat i=1
	=\lVert v \rVert \lVert \hat i \rVert \cos\theta
	=\sqrt{n}\cos\theta
	\implies \cos\theta=\frac{1}{\sqrt{n}}
\end{equation*}
As $n\to\infty$, $\sqrt{n}\to \infty$ and $1/ \sqrt{n}\to0$.
\begin{equation*}
    \cos\theta=0
	\implies \theta=\frac{\pi}{2}=\boxed{90^\circ}
\end{equation*}
The angle between the $x$-axis and the vector to one vertex of the hypercube tends to $90^\circ$ as the number of dimensions $n\to\infty$.

\problem{5}
Let $\hat u$ and $\hat v$ be unit vectors in the same direction as $u$ and $v$, respectively.
\begin{equation*}
    \hat u=\frac{u}{\lVert u \rVert } \qquad
    \hat v=\frac{v}{\lVert v \rVert } \qquad
	w'=\hat u+\hat v
\end{equation*}
Since $\hat u$ and $u$ point in the same direction, $w'$ bisects $u$ and $v$ if and only if it bisects $\hat u$ and $\hat v$. Let $\theta\in[0,\pi]$ be angle between $\hat u$ and $w'$, and let $\varphi\in[0,\pi]$ be the angle between $\hat v$ and $w'$.
\begin{alignat*}{3}
	\hat u\cdot w'
	=\lVert \hat u \rVert \lVert w' \rVert \cos\theta
	&\iff \hat u\cdot(\hat u+\hat v)
	=\lVert w' \rVert \cos\theta
	&\iff \cos\theta
	=\frac{\hat u\cdot \hat u+\hat u\cdot\hat v}{\lVert w' \rVert }
	=\frac{1+\hat u\cdot\hat v}{\lVert w' \rVert } \\
	\hat v\cdot w'
	=\lVert \hat v \rVert \lVert w' \rVert \cos\varphi
	&\iff \hat v\cdot(\hat u+\hat v)
	=\lVert w' \rVert \cos\varphi
	&\iff \cos\varphi
	=\frac{\hat v\cdot \hat u+\hat v\cdot\hat v}{\lVert w' \rVert }
	=\frac{1+\hat u\cdot\hat v}{\lVert w' \rVert }
\end{alignat*}
Cosine is injective on $[0,\pi]$ so $\cos\theta=\cos\varphi\implies \theta=\varphi$. Since the angle between $u$ and $w'$ is the same as the angle between $v$ and $w'$, the vector $w'=\hat u+\hat v$ bisects the angle between $u$ and $v$.

Notice that:
\begin{equation*}
    \lVert u \rVert \lVert v \rVert w'
	=\lVert u \rVert \lVert v \rVert \left( \frac{u}{\lVert u \rVert}+\frac{v}{\lVert v \rVert } \right)
	=\lVert v \rVert u+\lVert u \rVert v
	=w
\end{equation*}
Since $\lVert u \rVert \lVert v \rVert $ is a scalar, $w$ and $w'$ point in the same direction. $w$ bisects the angle between $u$ and $v$. \qed

\end{document}
